\section{Introduction}\label{sec:intro}

Mixed models with subject-level random effects are fitted by Hamiltonian
Monte Carlo (HMC) and its adaptive variant, the No-U-Turn sampler
\citep[NUTS;][]{hoffman2014}, in general-purpose software such as Stan
\citep{carpenter2017}, brms \citep{burkner2017}, rstanarm and NumPyro
\citep{phan2019}. brms and rstanarm generate the non-centred
parameterization of the random effects by default, and it is the standard
recommendation for models written directly in Stan or NumPyro: $b =
\Bstd L^{\top}$ with $\Bstd$ standard normal and $L$ a Cholesky factor of
the random-effect covariance. This
removes the funnel between the random effects and their scale
\citep{neal2003,betancourt2015} and is a main reason these samplers work well on hierarchical models.

It leaves a second, simpler problem in place. Suppose the fixed-effects
design contains a column that is constant within subject: the intercept,
a treatment arm, sex. Adding a constant $c$ to every subject's random
intercept and subtracting $c$ from the fixed intercept leaves every fitted
value unchanged. The likelihood identifies the sum; only the priors pin the
split, at scale $\sigma_0/\sqrt{N}$ for $N$ subjects when the fixed
intercept's prior is diffuse. The same holds for every subject-constant covariate, whose
coefficient trades off against a covariate-weighted combination of the
random intercepts, and, whenever there is a random slope on time, for the
fixed time slope and the mean random slope. In HMC's sampling coordinates
these directions are spread across all $N$ subjects, the posterior has a
ridge along them, and the sampler can mix slowly on the intercept, the
subject-constant coefficients and the time slope: the coefficients an
analyst reports. In a joint model of \texttt{pbc2}, for instance, the
posterior correlation between the fixed intercept and the mean random
intercept was $-0.97$ (Section~\ref{sec:degeneracy}).

The problem is old and well understood for Gibbs samplers. Hierarchical
centring \citep{gelfand1995} removes it by sampling the centred effects,
at the price of the funnel that non-centring was introduced to avoid, and
the choice between the two is the subject of a large literature
\citep{papaspiliopoulos2003,papaspiliopoulos2007,yu2011}. A second route
removes the directions: sum-to-zero constraints, whose effect on Gibbs and
on NUTS \citet{zanella2021} quantify; Stan's \texttt{sum\_to\_zero\_vector};
restricted spatial regression \citep{reich2006,hodges2010}; and, exactly,
the transformation of \citet{vines1996}, which integrates the grand-mean
directions out. The covariate-weighted direction for a random intercept
appears in the discussion of \citet{zanella2021} by \citet{yang2021} as a
device for analysing Gibbs samplers, and \citet{browne2009} name the
problem directly, suggesting as future work that predictors be made
orthogonal to the cluster indicators. Section~\ref{sec:related} places the
present work among these.

This paper takes a different route. It keeps the absorbable directions,
and makes them sampling coordinates. We construct an orthonormal basis for
every direction the fixed effects can absorb, in every random-effect
column, and apply one fixed rotation to the standardized random effects so
that these directions become a small block $U$ of $k \times q$ coordinates,
$k$ being the number of directions and $q$ the number of random-effect
columns. A single dense block of the HMC mass matrix over $(\beta, U)$
then largely whitens the ridge. The rotation acts on the subject index and the
covariance factor on the column index, so they commute: the construction is
exact for every value of the random-effect covariance, with no constraint,
no correction and no change to any reported quantity.

The motivating application is the Bayesian joint model of a longitudinal
marker and a time-to-event outcome
\citep{wulfsohn1997,henderson2000,tsiatis2004,rizopoulos2012}, as
implemented by the R package \texttt{jmjax}, which samples with NUTS in
NumPyro. There the random effects also enter the hazard through the current
marker value, and the survival submodel adds a condition on which
directions remain absorbable. Joint models also raise a second question.
Their primary estimand is the association parameter $\alpha$, the log
hazard ratio per unit of the current marker value. On it, NUTS gave 3 to 8 times the ESS per second of the Metropolis-within-Gibbs sampler of JMbayes2 \citep{jmbayes2}, with no reparameterization at all (3.5 and 8.2 times on two real datasets with one ESS estimator applied to both). We explain most of the per-draw part of that gap: for an ideal two-block sampler that draws $\alpha$ exactly given the rest, the
lag-1 autocorrelation equals the share of $\alpha$'s posterior variance
the rest explains \citep{liu1994}, and in the two joint models we measured that share was large, mostly through the survival parameters, plausibly because $\alpha$ trades off against the level of the log baseline hazard.

\paragraph{Contributions.}
\begin{enumerate}
\item \emph{The absorbable directions.} An exact characterization of the
random-effect directions fixed effects can absorb, for any design and every
random-effect column (Proposition~\ref{prop:char}); a construction that
depends only on the column space of the design, with an example where the
natural column-by-column check silently misses the slope direction
(Example~\ref{ex:spline}); and, for joint models, the condition the
survival submodel imposes (Condition~\ref{cond:S},
Proposition~\ref{prop:joint}).
\item \emph{An exact rotation and a small dense metric block.} One
Householder rotation of the standardized random effects, costing $O(Nkq)$
per gradient, localizes every absorbable direction in a $(p + kq)$-dimensional
block (Proposition~\ref{prop:exact}, Lemma~\ref{lem:local}). An analysis of
the resulting ridge (Proposition~\ref{prop:ridge}) explains why a diagonal
metric cannot handle it, why one dense block can, and where a fixed metric
is weakest; a stress grid aimed at those weak spots kept the ESS-per-second gains on the intercept and time slope at 5--30 times.
\item \emph{Blocked samplers and the association parameter.} The
application of the data-augmentation autocorrelation identity to
$\alpha$ (Proposition~\ref{prop:blocked}), with a bound that does not
need Gaussianity, the measurement of the relevant $R^2$ on real data
(0.85--0.87), and its decomposition, which
places most of it in the survival block rather than the random effects.
\item \emph{Evidence.} Simulation grids for one and two random-effect
columns, a 100-replicate calibration study, two linear and generalized
linear mixed models outside the joint-model setting, and two standard
joint-model datasets, including a comparison with JMbayes2.
\end{enumerate}

Section~\ref{sec:model} sets out the model and the degeneracy,
Section~\ref{sec:absorbable} the absorbable directions,
Section~\ref{sec:rotation} the rotation and the metric, and
Section~\ref{sec:blocked} the blocked-sampler result. Sections
\ref{sec:simulation}--\ref{sec:realdata} give the evidence and
Section~\ref{sec:discussion} discusses related work, costs and extensions.
Proofs not given in the text are in Appendix~\ref{app:proofs}; the scripts
that reproduce every table and figure are listed in
Appendix~\ref{app:repro}.
